\documentclass[preview]{standalone}

\usepackage{amsmath}
\usepackage{amssymb}
\usepackage{stellar}
\usepackage{definitions}

\begin{document}

\id{psicologia-apprendimento-cognitivo-situato}
\genpage

\section{Apprendimento cognitivo e situato}

\begin{snippetdefinition}{imprinting-definition}{Imprinting}
    L'\emph{imprinting} è una forma di apprendimento precoce osservabile in
    animali appena nati che, come i pulcini, mostrano una reazione di
    inseguimento verso il primo oggetto mobile che vedono o sentono.
\end{snippetdefinition}

\begin{snippetexample}{lorenz-imprinting-example}{Lorenz e l'imprinting}
    Konrad Lorenz mostrò che alcune oche appena nate potevano riconoscerlo come
    figura di riferimento e seguirlo come se fosse la madre. Interpretò questo
    fenomeno come una forma di apprendimento innato efficace solo entro un breve
    periodo iniziale dello sviluppo.
\end{snippetexample}

\includesnpt[width=55\%,src=/snippet/static-psychology/lorenz-imprinting-geese.jpg]{centered-img}

\begin{snippetdefinition}{periodo-critico-definition}{Periodo critico}
    Il \emph{periodo critico} è un periodo breve e circoscritto delle prime fasi
    dello sviluppo durante il quale alcune forme di apprendimento avverrebbero
    in modo particolarmente efficace.
\end{snippetdefinition}

\begin{snippetdefinition}{periodo-sensibile-definition}{Periodo sensibile}
    Il \emph{periodo sensibile} è il periodo in cui le influenze ambientali sono
    particolarmente efficaci per l'apprendimento di conoscenze e abilità, senza
    implicare necessariamente un limite rigido e irreversibile.
\end{snippetdefinition}

\begin{snippet}{periodo-sensibile-critica}
    Studi successivi, come quelli di Vallortigara e collaboratori, hanno
    criticato l'idea che l'imprinting sia sempre rapido e irreversibile. Per
    questo il concetto di periodo critico viene spesso sostituito con quello di
    periodo sensibile, applicato anche allo sviluppo del legame di attaccamento
    nella specie umana.
\end{snippet}

\begin{snippetdefinition}{apprendimento-cognitivo-definition}{Apprendimento cognitivo}
    L'\emph{apprendimento cognitivo} è una forma di apprendimento spiegata in
    termini di cambiamenti nei processi mentali, piuttosto che soltanto in
    termini di cambiamenti comportamentali osservabili.
\end{snippetdefinition}

\begin{snippet}{apprendimento-cognitivo-contesto}
    Alcune forme di apprendimento sono difficili da spiegare nei termini del
    condizionamento classico o operante. Secondo la psicologia cognitiva, esse
    richiedono di considerare i cambiamenti nelle rappresentazioni mentali e nel
    modo in cui l'individuo organizza il problema.
\end{snippet}

\begin{snippetdefinition}{apprendimento-insight-definition}{Apprendimento per insight}
    L'\emph{apprendimento per insight} è una forma di apprendimento che avviene
    attraverso una riorganizzazione improvvisa della percezione del problema in
    una data situazione.
\end{snippetdefinition}

\begin{snippetexample}{sultan-kohler-example}{Sultan}
    Wolfgang Köhler studiò lo scimpanzé Sultan, che aveva già imparato ad
    ammucchiare scatole per raggiungere frutta appesa e a usare un bastone per
    prendere frutta fuori dalla sua portata. In una situazione nuova, Sultan
    fallì inizialmente per prove ed errori; dopo una pausa, combinò i due
    apprendimenti, trascinò una cassa e un bastone sotto la frutta, salì sulla
    cassa e raggiunse il cibo con il bastone. Il comportamento fu interpretato
    come apprendimento per insight.
\end{snippetexample}

\includesnpt[width=58\%,src=/snippet/static-psychology/kohler-sultan-insight-learning.jpg]{centered-img}

\begin{snippetdefinition}{mappa-cognitiva-definition}{Mappa cognitiva}
    Una \emph{mappa cognitiva} è una rappresentazione mentale che un organismo
    utilizza per orientarsi e muoversi all'interno di un ambiente familiare.
\end{snippetdefinition}

\begin{snippetexample}{tolman-mappe-cognitive-example}{Tolman e le mappe cognitive}
    Tolman osservò che i ratti potevano scegliere rapidamente percorsi
    alternativi in un labirinto quando la strada abituale verso l'obiettivo
    veniva bloccata. Questo comportamento suggeriva che non stessero eseguendo
    solo una catena meccanica di risposte, ma usassero una rappresentazione
    mentale dell'ambiente.
\end{snippetexample}

\includesnpt[width=48\%,src=/snippet/static-psychology/tolman-cognitive-map-maze.jpg]{centered-img}

\begin{snippetdefinition}{apprendimento-latente-definition}{Apprendimento latente}
    L'\emph{apprendimento latente} è una forma di apprendimento che avviene
    senza manifestarsi immediatamente nel comportamento e senza rinforzi
    evidenti, ma che può emergere in seguito quando diventa utile.
\end{snippetdefinition}

\begin{snippetdefinition}{apprendimento-situato-definition}{Apprendimento situato}
    L'\emph{apprendimento situato} è un apprendimento legato a una specifica
    situazione e immerso in un determinato contesto fisico, sociale e relazionale.
\end{snippetdefinition}

\begin{snippet}{apprendimento-situato-caratteri}
    Nell'apprendimento situato, esperto e novizio sono interdipendenti nella
    costruzione degli apprendimenti. L'apprendimento è interattivo, aperto,
    individualizzato e contingente, perché dipende dalle opportunità e dai
    vincoli del contesto in cui avviene.
\end{snippet}

\begin{snippetdefinition}{conoscenza-tacita-definition}{Conoscenza tacita}
    La \emph{conoscenza tacita} è una conoscenza individuale implicita,
    difficilmente codificabile, trasmissibile o condivisibile in modo esplicito.
    È in larga misura inconsapevole e immersa nell'esperienza pratica.
\end{snippetdefinition}

\begin{snippetdefinition}{e-learning-definition}{E-learning}
    L'\emph{e-learning} è una forma di apprendimento che usa la tecnologia come
    strumento privilegiato.
\end{snippetdefinition}

\begin{snippetdefinition}{serious-games-definition}{Serious Games}
    I \emph{Serious Games} sono attività digitali interattive che, attraverso la
    simulazione virtuale, permettono ai partecipanti di fare esperienze precise
    e accurate, promuovendo percorsi attivi, partecipati e coinvolgenti di
    apprendimento.
\end{snippetdefinition}

\begin{snippet}{serious-games-elementi}
    I Serious Games collegano gioco, simulazione mentale e apprendimento. Di
    solito includono un obiettivo sfidante, una struttura coinvolgente,
    l'attribuzione di un punteggio e la possibilità di acquisire abilità,
    atteggiamenti o conoscenze applicabili al mondo reale.
\end{snippet}

\includesnpt[width=58\%,src=/snippet/static-psychology/serious-games-learning.jpg]{centered-img}

\begin{snippet}{valutazione-apprendimento}
    La valutazione è un momento essenziale dell'apprendimento, perché permette
    di verificare quanto il discente abbia imparato e in che modo. Gli strumenti
    tradizionali, come interrogazioni, colloqui e test, devono rispettare criteri
    di attendibilità, validità, accuratezza e finezza discriminativa. I Serious
    Games possono integrare questi strumenti offrendo valutazioni dinamiche e
    feedback in tempo reale.
\end{snippet}

\end{document}
